\section{Curves realizing binary power addition}\label{sec:curves}

We next construct commuting matrix polynomials that replace two inputs by
one output whose \(b\)-th power is their sum.  The matrices vary over a
Picard variety.  An accompanying finite group action, free on a
degree-one Picard variety, will allow the next two sections to control
the division algebra obtained by descent.

All curves in this section are defined over the original field \(\C\).
They may therefore be extended to any of the centers used in the
construction.  For a left group action, transport of line bundles and
sections means \(g_*=(g^{-1})^*\).  In particular the action on
\(\Pic^d(C)\) is \([A]\mapsto[g_*A]\).
A \(G\)-linearization of a line bundle means lifts of the action to its
total space that are linear on fibers and satisfy the group law exactly.
This is stronger than isomorphisms between transported bundles whose
composites agree only up to a scalar.

\begin{proposition}[Binary addition data]\label{prop:binary-curves}
Let \(b\ge2\), fix a primitive \(b\)-th root \(\zeta\), and put
\[
 G=(\mu_b^b/\mu_b)\rtimes(\Z/b),\qquad e=b^b,
\]
where the denominator is the diagonal subgroup and \(\Z/b\) permutes
the \(b\) factors cyclically.  There are smooth connected projective
curves \(C,T_0\), a connected \'etale \(G\)-cover \(C\to T_0\), and a
genuinely \(G\)-linearized line bundle \(L\) on \(C\) with the following
properties.
\begin{enumerate}[label=\textup{(\roman*)}]
\item The genera and degree satisfy
\[
 g_T:=g(T_0)=b-1,\qquad g(C)-1=e(g_T-1),\qquad \deg L=e.
\]
\item There are sections \(a_\alpha\in H^0(C,L)\), indexed by
\(\alpha\in\Z_{\ge0}^b\) with \(|\alpha|=b\), whose transport follows
the degree-\(b\) monomial representation in formal homogeneous
coordinates \(p_0,\ldots,p_{b-1}\). The notation \(p^\alpha\) labels
the section \(a_\alpha\) when \(|\alpha|=b\); it does not assign
individual sections of \(L\) to the \(p_j\). Write \(a_\alpha=p^\alpha\)
and set
\[
 z=\frac1b\sum_{j=0}^{b-1}p_j^b,\qquad
 u=-\frac1b\sum_{j=0}^{b-1}\zeta^{-j}p_j^b,\qquad
 v=p_0\cdots p_{b-1}.
\]
They satisfy
\begin{equation}\label{eq:binary-section-relations}
 p_j^b=z-\zeta^ju,\qquad z^b=u^b+v^b.
\end{equation}
The section \(z\) is invariant, and \(u,v\) have characters
\(\chi,\psi:G\to\mu_b\), respectively.  The pencil \(u,v\) is
base-point-free and defines a degree-\(e\) morphism
\(\pi:C\to\PP^1\) with \(L=\pi^*\cO(1)\).
\item For a generic \(A\in\Pic^0(C)\), over its field of definition
\(\kappa\), the space \(E_A=H^0(C,A L^{g_T})\) has dimension \(e\).
Multiplication of sections gives isomorphisms
\begin{equation}\label{eq:curve-free-module}
 H^0(C,A L^{g_T+j})
   =E_A\otimes_\kappa\Sym^j_\kappa\langle u,v\rangle
       \qquad(j\ge0).
\end{equation}
Every \(a\in H^0(C,L)\) consequently acts by a homogeneous linear
matrix polynomial \(M_a(U,V)\in\End_\kappa(E_A)[U,V]\).
These matrix polynomials commute, preserve all polynomial identities
among the sections, and satisfy
\[
 M_u=U\,\mathrm{id},\qquad M_v=V\,\mathrm{id},\qquad
 M_z^b=(U^b+V^b)\,\mathrm{id}.
\]
Their dependence on \(A\) is algebraic on a nonempty \(G\)-stable
open subset of \(\Pic^0(C)\).  Under transport on this family,
\begin{equation}\label{eq:matrix-transport}
 M_{g_*a}(U,V)
     =(g_{\mathrm{coeff}}M_a)(\chi(g)U,\psi(g)V).
\end{equation}
Here \(g_{\mathrm{coeff}}\) transports the endomorphism coefficients;
scalar choices in identifying transported line bundles have no effect.
\item The action of \(G\) on \(\Pic^1(C)\) is free.  For every character
\(\xi:G\to\mu_b\), there is a continuous map
\(\Pic^1(C)\to S^1\) equivariant for the given \(G\)-action and
multiplication by \(\xi\) on \(S^1\).
\end{enumerate}
\end{proposition}

The three sections \(u,v,z\) give the binary identity. The full monomial
family \(a_\alpha\) also encodes the construction by linear and quadratic
relations in Section~\ref{sec:assembly}, as required by the induction
data of Section~\ref{sec:low-relations}.
We prove the proposition in three steps: a cover with the required
sections, a free module for multiplication by those sections, and the
two properties of the degree-one Picard variety.

\subsection{The cover and its sections}

\begin{lemma}\label{lem:binary-cover}
There are curves, a cover, and a linearized line bundle satisfying
parts \textup{(i)} and \textup{(ii)} of
Proposition~\ref{prop:binary-curves}.
\end{lemma}

\begin{proof}
Inside \(\PP^{b-1}\), with coordinates \(p_0,\ldots,p_{b-1}\), take
the inverse image under coordinate-wise \(b\)-th powering of the line
parameterized by
\[
 [u:z]\longmapsto[z-u:z-\zeta u:\cdots:z-\zeta^{b-1}u].
\]
Thus the preimage is defined by the \(b-2\) independent linear
relations on the \(p_j^b\); no equation is imposed when \(b=2\).
Let \(C'\) be its normalization.  To establish its irreducibility,
put \(s=u/z\).  On the generic fiber over the \(s\)-line the coordinate
ratios satisfy
\[
 (p_j/p_0)^b=\frac{1-\zeta^j s}{1-s},\qquad 1\le j<b.
\]
These \(b-1\) functions generate a free
\((\Z/b)\)-submodule of
\(\C(s)^\times/\C(s)^{\times b}\): valuation at
\(s=\zeta^{-j}\) detects the \(j\)-th exponent independently.
Kummer theory gives a field extension of degree \(b^{b-1}\).
This argument uses valuations equal to one and applies also when
\(b\) is composite.  Every component of the projective preimage
has dimension at least one, since there are \(b-2\) equations, and
the coordinate-power map is finite.  Hence every component dominates
the line.  The integral generic fiber therefore proves that its
reduced support, and consequently \(C'\), is irreducible.

Phase changes of the \(p_j\), together with their cyclic permutation,
act on this curve.  The group before projectivization is
\(\widetilde G=\mu_b^b\rtimes(\Z/b)\).  Its scalar subgroup is the
diagonal \(\mu_b\), and the cyclic permutation already has order
\(b\) before taking the quotient.  Thus the resulting group is
the stated semidirect product \(G\).  Its action is faithful:
a nontrivial cyclic permutation moves \(s\) by a nontrivial
\(b\)-th root, whereas a phase change fixing a general point with
all coordinates nonzero must be diagonal.
The invariant coordinate
\[
 t=s^b
\]
defines a degree-\(b^b\) map \(C'\to\PP^1_t\).  Since \(|G|=b^b\),
this is a Galois cover with group \(G\).

The Kummer cover of the \(s\)-line is ramified precisely at the
\(b\) points \(s=\zeta^{-j}\), with index \(b\).  In particular,
although several of the displayed ratios have a pole at \(s=1\),
their simultaneous adjunction still has local index \(b\): every
unit in \(\C[[x]]\) has a \(b\)-th root, so all local extensions
are contained in \(\C((x^{1/b}))\).
The map is unramified over \(s=0,\infty\).
It follows that the map to the \(t\)-line has exactly the branch
values \(0,1,\infty\), with index \(b\) above each.

Let \(T_0\) be the smooth projective model of
\begin{equation}\label{eq:auxiliary-cyclic-cover}
 d^b=\frac{t(t-1)}{t-2}.
\end{equation}
The valuations of the right side at \(0,1,2,\infty\) are
\(1,1,-1,-1\), respectively.  This is a cyclic degree-\(b\)
cover, totally ramified at those four points and unramified
elsewhere.  Each nontrivial intermediate extension is ramified
at \(2\), whereas \(C'\to\PP^1_t\) is unramified there.
The two Galois extensions are consequently linearly disjoint.
Normalize their fiber product to obtain \(C\).
It is connected, and \(C\to T_0\) remains a \(G\)-cover.

This last cover is \'etale.  Above \(0,1,\infty\), both original
ramification indices are \(b\); in completed local coordinates,
adjoining the same \(b\)-th root of a uniformizer removes the
ramification after normalization.  Above \(2\), and away from
the four listed points, the base change of the original map is
already \'etale.  Riemann--Hurwitz now gives
\[
 2g_T-2=-2b+4(b-1)=2b-4,\qquad
 g(C)-1=|G|(g_T-1).
\]
For \(b=2\), this construction starts with \(C'=\PP^1\) and
produces a degree-four \'etale \(C_2\times C_2\)-cover between
curves of genus one.

Pull \(\cO(b)|_{C'}\) back to \(C\), and call the resulting
bundle \(L\).  The sections \(u,z\) defined in the proposition
have no common zero, since their simultaneous vanishing forces
every \(p_j\) to vanish.  On \(C'\) their pencil is the
degree-\(b^{b-1}\) map to \(\PP^1_s\); hence
\(\deg\cO(b)|_{C'}=b^{b-1}\).  Since \(C\to C'\) has degree
\(b\), we have \(\deg L=b^b=e\).

The standard \(\widetilde G\)-linearization on \(\cO(b)\)
descends to a genuine \(G\)-linearization: a diagonal scalar
in \(\mu_b\) acts trivially in degree \(b\).  Pullback preserves
this linearization.  Degree-\(b\) monomials therefore define
the sections \(a_\alpha\) and their asserted action.
Phase changes fix \(u,z\), cyclic permutation fixes \(z\)
and multiplies \(u\) by a primitive root, and \(v=\prod_jp_j\)
transforms by the product of phases.  This product is invariant
under cyclic permutation and trivial on the diagonal subgroup,
so it defines a character of \(G\).  Denote the resulting
transport characters of \(u,v\) by \(\chi,\psi\).

The equations of the curve give the first relations in
\eqref{eq:binary-section-relations}, and their product gives
\[
 v^b=\prod_{j=0}^{b-1}(z-\zeta^ju)=z^b-u^b.
\]
If \(u=v=0\), this identity gives \(z=0\), which was excluded.
Thus \(u,v\) define a base-point-free pencil.  Its morphism
\(\pi:C\to\PP^1\) satisfies \(L=\pi^*\cO(1)\); since
\(\deg L=e>0\), the morphism is finite of degree \(e\).
\end{proof}

Figure~\ref{fig:curve-covers} summarizes the normalization and its two
projections. It separates the original power-cover geometry from the
base change that removes its ramification.

\begin{figure}[!htbp]
\centering
\begin{tikzpicture}[>=Stealth,every node/.style={font=\small}]
\node (C) at (0,2) {$C$};
\node (Cp) at (5,2) {$C'$};
\node (T) at (0,0) {$T_0$};
\node (P) at (5,0) {$\PP^1_t$};
\draw[->] (C)--node[above]{degree $b$}(Cp);
\draw[->] (C)--node[left]{\'etale $G$-cover}(T);
\draw[->] (Cp)--node[right]{degree $e=b^b$}(P);
\draw[->] (T)--node[below]{degree $b$}(P);
\end{tikzpicture}
\caption{The curve $C$ is the normalization of the fiber product.
The lower cover is totally ramified at $0,1,2,\infty$; its ramification
at $0,1,\infty$ cancels that of the right-hand cover, making the
left-hand cover \'etale.}
\label{fig:curve-covers}
\end{figure}
\FloatBarrier

\subsection{Multiplication as linear matrix polynomials}

The cover supplies the required power identity among sections.
We now turn those sections into linear matrix polynomials, with
no choice of basis built into the construction. The use of a trivial
finite pushforward to linearize multiplication is the vector-bundle
description of matrix representations of power identities
\cite{VanDenBergh87,Kulkarni03,CKM12}; we give the needed argument
for this curve and pencil explicitly.

\begin{lemma}\label{lem:binary-matrices}
For the curve and line bundle of Lemma~\ref{lem:binary-cover},
part \textup{(iii)} of Proposition~\ref{prop:binary-curves} holds.
\end{lemma}

\begin{proof}
Let \(g=g(C)\).  For generic \(A\in\Pic^0(C)\), the line bundle
\(A L^{g_T-1}\) is generic of degree
\[
 e(g_T-1)=g-1.
\]
The effective locus in \(\Pic^{g-1}(C)\) is the image of
\(\Sym^{g-1}C\), a projective variety of dimension \(g-1\).
It is a proper closed subset of the \(g\)-dimensional Picard
variety.  Thus \(H^0(C,A L^{g_T-1})=0\), and
Riemann--Roch gives \(H^1(C,A L^{g_T-1})=0\) as well.
Here \(g\ge1\); when \(g=1\), the effective locus in degree
zero consists only of the trivial line bundle.
The complement of this locus, translated to \(\Pic^0(C)\),
is a nonempty \(G\)-stable open subset, because \(L\) is
linearized.

Fix \(A\) in that open set, and put
\[
 \mathcal V=\pi_*(A L^{g_T}).
\]
This is a rank-\(e\) vector bundle on \(\PP^1\): the map
\(\pi\) is finite, and its pushforward is torsion-free on a
smooth curve.  Projection formula and the preceding
vanishing give \(H^*(\PP^1,\mathcal V(-1))=0\).
For every geometric point \(q\in\PP^1\), the exact sequence
\[
 0\longrightarrow\mathcal V(-q)\longrightarrow\mathcal V
   \longrightarrow\mathcal V|_q\longrightarrow0
\]
therefore identifies \(H^0(\PP^1,\mathcal V)\) with
\(\mathcal V|_q\).  It follows that the evaluation map
\[
 H^0(\PP^1,\mathcal V)\otimes\cO_{\PP^1}
       \longrightarrow\mathcal V
\]
is an isomorphism.  This also proves the assertion over the
field of definition of \(A\), since an isomorphism can be
checked after algebraic closure.  In particular
\(\dim E_A=e\).  Tensoring by \(\cO(j)\), taking global
sections, and using projection formula proves
\eqref{eq:curve-free-module}, with its indicated multiplication
map.

Let
\[
 \mathcal M_A=\bigoplus_{j\ge0}H^0(C,A L^{g_T+j}).
\]
Equation~\eqref{eq:curve-free-module} identifies this graded
module canonically with \(E_A\otimes_\kappa\kappa[U,V]\).
Multiplication by a section \(a\) of \(L\) is a homogeneous
degree-one endomorphism of this free module over
\(\kappa[U,V]\).  Such an endomorphism is determined by its
value on degree zero and has the form
\[
 M_a(U,V)=A_aU+B_aV,\qquad A_a,B_a\in\End_\kappa(E_A).
\]
Multiplication of sections is commutative, so these matrix
polynomials commute as elements of \(\End_\kappa(E_A)[U,V]\).
Every polynomial identity among sections acts as the zero
endomorphism and remains an identity among the matrices.
In particular \(u,v\) act as \(U,V\), and the identity for
\(M_z\) follows from \eqref{eq:binary-section-relations}.
This does not assert that all the individual matrices
\(A_a,B_a\) commute.

Choose a point of the constant curve to normalize the
universal line bundle on \(\Pic^0(C)\times C\).
On the open subset just specified, the dimensions of the
section spaces are constant and the higher cohomology
vanishes in every degree \(g_T+j\), \(j\ge0\).
Cohomology and base change, applied to these line bundles along the
proper constant-curve projection, then makes the preceding evaluation
and multiplication maps algebraic in \(A\)
\cite[Lemma~30.22.1, Tag 07VK]{Stacks}. The universal lines are flat
over the Noetherian parameter base, as required there.

Finally, all multiplication maps are natural under transport.
An isomorphism between a transported universal line and
the universal line at the transported parameter is unique
up to a scalar.  It acts by the same scalar on the source
and target section spaces, so the induced transport on
endomorphisms is independent of that choice and composes
as a group action.  Transporting
\[
 m_a:E_A\longrightarrow E_A\otimes\langle u,v\rangle
\]
and using \(g_*u=\chi(g)u\), \(g_*v=\psi(g)v\) gives
\eqref{eq:matrix-transport}.  On the generic parameter
field this is semilinear transport of the coefficients.
\end{proof}

\subsection{The degree-one Picard variety}

The multiplication construction uses \(\Pic^0(C)\).
The next two properties of \(\Pic^1(C)\) provide, respectively,
a free algebraic action for descent and continuous maps
carrying its prescribed characters.

\begin{lemma}\label{lem:picard-free}
If a finite group \(G\) acts freely on a smooth connected
projective complex curve \(C\), then its induced action on
\(\Pic^1(C)\) is free.
\end{lemma}

\begin{proof}
Suppose an element \(g\ne1\), of order \(r\), fixes
\([V]\in\Pic^1(C)\).  Choose an isomorphism \(g_*V\to V\).
Its \(r\)-fold composite is an automorphism of \(V\), hence
a nonzero scalar.  Rescaling the chosen isomorphism by
an \(r\)-th root of the inverse scalar gives a linearization
of \(V\) under \(\langle g\rangle\).
Since this cyclic subgroup acts freely on \(C\), descent
along the degree-\(r\) \'etale cover
\(q:C\to C/\langle g\rangle\) gives \(V=q^*V_0\).
But then \(1=\deg V=r\deg V_0\), a contradiction.
\end{proof}

\begin{lemma}\label{lem:picard-character-maps}
Let \(q:C\to T\) be a connected \'etale \(G\)-cover of smooth
connected projective complex curves.  For every character
\(\xi:G\to\mu_b\), there is a continuous map
\[
 f_\xi:\Pic^1(C)\longrightarrow S^1,\qquad
 f_\xi(g_*A)=\xi(g)f_\xi(A).
\]
\end{lemma}

\begin{proof}
Identify \(\mu_b\) with \(\Z/b\) by exponentiation.
Compose the monodromy of the cover with \(\xi\) to obtain
\(\bar\alpha\in H^1(T;\Z/b)\).
Since \(H_1(T;\Z)\) is free abelian, \(\bar\alpha\) lifts
to a class \(\alpha\in H^1(T;\Z)\).
Represent \(\alpha\) by a harmonic real one-form with
integral periods.  The form
\[
 \omega=\frac1b q^*\alpha
\]
has integral periods on \(C\): the projection of a closed
loop in \(C\) has trivial covering monodromy, so its
\(\alpha\)-period is divisible by \(b\).
It is also harmonic, since \(q\) is holomorphic and
unramified.  Integrating \(\omega\) modulo \(\Z\) and
exponentiating defines a circle-valued map on \(C\).
For a path from \(c\) to \(gc\), the projected loop has
character \(\xi(g)\); choosing the sign of the monodromy
convention accordingly, the map satisfies
\[
 f_C(gc)=\xi(g)f_C(c).
\]
Its differential, expressed as a real one-form after local
lifting to \(\R\), is the \(G\)-invariant form \(\omega\).

The Abel map
\[
 a:C\longrightarrow\Pic^1(C),\qquad c\longmapsto\cO(c),
\]
is equivariant.  By the Jacobian description as the
quotient of the dual space of holomorphic differentials
by the integral period lattice, pullback by \(a\)
identifies translation-invariant real one-forms on
\(\Pic^1(C)\) with harmonic real one-forms on \(C\),
and induces an isomorphism on integral first cohomology;
see \cite[Proposition~2.2 and Theorem~2.5]{MilneJacobian}.
Consequently \(\omega=a^*\widetilde\omega\), where
\(\widetilde\omega\) is translation-invariant with integral
periods.  Integrating \(\widetilde\omega\) gives a continuous
circle-valued map \(f_\xi\) on \(\Pic^1(C)\).  Adjust its
constant so that \(f_\xi\circ a=f_C\).

The action on the Picard torsor is affine.  The difference
between the pullback of \(\widetilde\omega\) under \(g\)
and \(\widetilde\omega\) is translation-invariant and pulls
back by \(a\) to zero, since \(\omega\) is invariant.
Injectivity of this pullback makes the difference zero.
It follows that \(f_\xi(g_*A)/f_\xi(A)\) is constant on the
connected Picard variety.  On the Abel image it equals
\(\xi(g)\), so it has that value everywhere.
\end{proof}

\begin{proof}[Proof of Proposition~\ref{prop:binary-curves}]
Lemma~\ref{lem:binary-cover} supplies parts \textup{(i)}
and \textup{(ii)}, and Lemma~\ref{lem:binary-matrices}
supplies part \textup{(iii)}.
Apply Lemmas~\ref{lem:picard-free}
and~\ref{lem:picard-character-maps} to the \'etale cover
\(C\to T_0\) to obtain part \textup{(iv)}.
\end{proof}
